% Figure: Murray's law branching geometry, with symmetric bifurcation
% angles and the cubed-radius constraint.
\begin{figure}[htbp]
\centering
\begin{tikzpicture}[
  vessel/.style={draw=engred, fill=engred!20, line width=2pt},
  arrow/.style={-{Stealth}, thick, draw=engblue},
  label/.style={font=\scriptsize, engblue},
]

% Parent vessel (vertical at the centre, large diameter)
\draw[vessel] (-0.45, 3.5) -- (-0.45, 1.8) -- (0.45, 1.8) -- (0.45, 3.5) -- cycle;

% Daughter vessels at +/- 37 degrees (symmetric, smaller diameter)
% Daughter half-width corresponds to (1/2)^(1/3) factor relative to parent
% so each daughter half-width ~ 0.36 (parent half-width was 0.45)
\begin{scope}[shift={(0, 1.8)}, rotate=-37]
  \fill[engred!20, draw=engred, line width=2pt]
    (-0.36, 0) -- (-0.36, -2.2) -- (0.36, -2.2) -- (0.36, 0) -- cycle;
\end{scope}
\begin{scope}[shift={(0, 1.8)}, rotate=37]
  \fill[engred!20, draw=engred, line width=2pt]
    (-0.36, 0) -- (-0.36, -2.2) -- (0.36, -2.2) -- (0.36, 0) -- cycle;
\end{scope}

% Angle arc
\draw[arrow] (0, 1.8) -- (0, 2.7);
\draw[arrow] (0, 1.8) -- (-37:1.6);
\draw[arrow] (0, 1.8) -- (37:1.6) coordinate (a37);
\draw[engblue, thick] (0.7, 2.4) arc (90:53:0.7);
\node[label] at (0.6, 2.0) {$\theta$};

% Vessel diameter labels
\node[label, anchor=west] at (0.7, 3.0) {$d_p$};
\node[label, anchor=west, rotate=-37] at (-37:-2.4) {$d_1$};
\node[label, anchor=west, rotate=37] at (37:-2.4) {$d_2$};

% Murray's law text
\node[font=\small] at (0, -2.5) {$d_p^3 = d_1^3 + d_2^3$};
\node[font=\small] at (0, -3.1) {symmetric: $d_1 = d_2$, so $d_1/d_p = 2^{-1/3} \approx 0.794$};
\node[font=\small] at (0, -3.7) {optimal $\theta \approx 37.5^{\circ}$ for symmetric bifurcation};

\node[font=\bfseries] at (0, 4.0) {Vascular bifurcation: Murray's law geometry};

\end{tikzpicture}
\caption{A vascular bifurcation under Murray's law (Murray 1926).
Minimising the sum of metabolic maintenance cost (proportional to
vessel volume) and Poiseuille pumping work (inversely proportional to
$d^4$) gives the cubed-radius constraint
$d_p^3 = d_1^3 + d_2^3$. For a symmetric bifurcation $d_1 = d_2$, each
daughter has diameter $2^{-1/3} \approx 0.794$ times the parent.
The optimal bifurcation angle is approximately $37.5^{\circ}$ for the
symmetric case. Both the radius ratio and the branching angle are
observed in mammalian arterial trees to within $\sim 10\,\%$ across
many orders of magnitude in vessel size, and they are the
engineering basis of scaffold designs that try to recapitulate
natural vasculature.}
\label{fig:vol06ch05:murray}
\end{figure}
